\section{Central Dedekind branches and an unconditional arithmetic imbalance}
\label{sec:analytic-central-puiseux}

\begin{theorem}[NN-A018: the second central Puiseux coefficient]
Let $K$ have degree $n=r_1+2r_2$ and absolute discriminant $D_K$.
Write $L=\zeta_K$ and suppose its zero at $\rho=1/2$ has order
$m\ge2$. Put
\[
 Q_K(s)=s(s-1)D_K^{s/2}\Gamma_{\mathbb R}(s)^{r_1}
           \Gamma_{\mathbb C}(s)^{r_2},\qquad
 \Gamma_{\mathbb R}(s)=\pi^{-s/2}\Gamma(s/2),\quad
 \Gamma_{\mathbb C}(s)=(2\pi)^{-s}\Gamma(s),
\]
so that $\xi_K=Q_KL$ is entire and satisfies
$\xi_K(s)=\xi_K(1-s)$. The harmless constant often inserted in
$\Gamma_{\mathbb C}$ changes none of the formulas below. Set
\[
 a_m=\frac{L^{(m)}(1/2)}{m!},\qquad
 C_K=2+\frac{\zeta_K'}{\zeta_K}(3/2)
                 +\frac{Q_K'}{Q_K}(1/2).
\]
The shifted Dirichlet-series deformation from NN-A004 has exactly $m$
local branches, counted with multiplicity. For positive real $t$ they
are distinct for all sufficiently small $t$ and have expansions
\begin{equation}\label{eq:central-second-jet}
 s_j(t)=\frac12+c_jt^{1/m}
              +\frac{C_K}{m}c_j^2t^{2/m}+O(t^{3/m}),\qquad
 c_j^m=\frac{\zeta_K(3/2)}{2a_m}.
\end{equation}
The central order $m$ is even. If $a_m<0$, let
$a=(\zeta_K(3/2)/(2|a_m|))^{1/m}$ and choose
$c_j=a\exp((2j+1)\pi i/m)$. If $m\equiv0\pmod4$, exactly $m/2$
branches lie on each side of the critical line for small positive $t$.
If $m\equiv2\pmod4$, precisely two leading directions are vertical.
When $C_K>0$ these two branches bend left, and the left/right counts
are $(m+2)/2,(m-2)/2$; when $C_K<0$ both bend right and these counts
are reversed. No conclusion about their normal direction is made
from this coefficient when $C_K=0$.

The hypothesis $a_m<0$ follows from GRH for $\zeta_K$, hence from A
for abelian $K$ and conditionally from M in general. It can also be
established without GRH by a local factorization, as in the example
below.
\end{theorem}

\begin{proof}
Evenness of $\xi_K(1/2+z)$ and the nonvanishing of $Q_K(1/2)$ imply
that $m$ is even. Write
\[
 L(1/2+z)=a_mz^m+a_{m+1}z^{m+1}+O(z^{m+2}).
\]
The coefficient of $z^{m+1}$ in $Q_K(1/2+z)L(1/2+z)$ vanishes, so
\begin{equation}\label{eq:central-unit-jet}
 \frac{a_{m+1}}{a_m}=-\frac{Q_K'}{Q_K}(1/2).
\end{equation}
The deformation has the local expansion
\[
 L_t(1/2+z)=a_mz^m+a_{m+1}z^{m+1}
       -t(B_0+B_1z)+O(z^{m+2})+O(tz^2)+O(t^2),
\]
where $B_0=L(3/2)/2>0$ and
$B_1=L(3/2)+L'(3/2)/2$. By NN-A005, the base change $t=u^m$
produces $m$ holomorphic branches $z=c_ju+d_ju^2+O(u^3)$,
with distinct $c_j$ and $a_mc_j^m=B_0$. The coefficient of $u^{m+1}$
is
\[
 ma_mc_j^{m-1}d_j+a_{m+1}c_j^{m+1}-B_1c_j=0.
\]
Because $2m\ge m+2$, the $O(t^2)$ term does not contribute.
Using $a_mc_j^m=B_0$ and \eqref{eq:central-unit-jet} gives
\[
 d_j=\frac{c_j^2}{m}
       \left(\frac{B_1}{B_0}-\frac{a_{m+1}}{a_m}\right)
     =\frac{C_K}{m}c_j^2.
\]
This proves \eqref{eq:central-second-jet} and the simplicity of
the branches at small nonzero $t$.

For $a_m<0$ the displayed $m$ directions are the roots of a negative
real number. A direction is vertical exactly when $m\equiv2\pmod4$,
and then the two vertical directions are $\pm ia$. Their second
normal coefficient is $-a^2C_K/m$, the same for both. The other
directions pair across the imaginary axis. This gives all the
claimed counts. The $O(t^{3/m})$ error is smaller than the nonzero
first or second normal term, so the statements hold on a common
positive interval for all $m$ branches.

Under GRH, $\xi_K$ has no real zero in $(1/2,\infty)$ and is positive
on $(1,\infty)$. Its value at $1$ is the positive Dedekind residue
times positive factors; thus it remains positive on $(1/2,\infty)$.
Its first nonzero coefficient at $1/2$ is consequently positive.
Since $Q_K(1/2)<0$, this gives $a_m<0$.
\end{proof}

\begin{proposition}[An effective conductor criterion]
Let
\[
 a_{\mathbb R}=\frac12\psi(1/4)-\frac12\log\pi,
 \qquad C_*=-a_{\mathbb R}-\frac{\zeta'}{\zeta}(3/2).
\]
Then
\begin{equation}\label{eq:central-C-bound}
 C_K\ge2+\frac12\log D_K-nC_*+\frac\pi2r_2.
\end{equation}
In particular, root discriminant $D_K^{1/n}\ge4371$ implies $C_K>0$.
Consequently, whenever $a_m<0$ and $m\equiv2\pmod4$, the two
initially vertical branches bend left for such a field.
\end{proposition}

\begin{proof}
The logarithmic derivative of $s(s-1)$ vanishes at $1/2$.
The gamma reflection and duplication formulas give
\[
 \frac{Q_K'}{Q_K}(1/2)=\frac12\log D_K
                   +na_{\mathbb R}+\frac\pi2r_2.
\]
For a rational prime $p$, write $f_{\mathfrak p}$ and
$e_{\mathfrak p}$ for the residue and ramification degrees.
For $\sigma>1$,
\[
 -\frac{\zeta_K'}{\zeta_K}(\sigma)
 =\sum_p\log p\sum_{\mathfrak p\mid p}
       \frac{f_{\mathfrak p}}{p^{f_{\mathfrak p}\sigma}-1}
 \le n\sum_p\frac{\log p}{p^\sigma-1}
 =-n\frac{\zeta'}{\zeta}(\sigma).
\]
Here $p^{f\sigma}-1\ge p^\sigma-1$ and
$\sum f_{\mathfrak p}\le\sum e_{\mathfrak p}f_{\mathfrak p}=n$.
This proves \eqref{eq:central-C-bound}. The saved Arb certificate
gives $e^{2C_*}<4371$ by a strict interval comparison; hence the
right side of \eqref{eq:central-C-bound} is positive at the stated
root discriminant. One may instead use the exact threshold
$\log D_K>2nC_*-\pi r_2-4$ without decimal constants.
\end{proof}

\begin{corollary}[An unconditional arithmetic example with two excess right branches]
Let
\[
 F=\mathbb Q(\sqrt5,\sqrt{41}),\qquad
 N=F\left(\sqrt{\frac{(5+\sqrt5)(41+\sqrt{205})}{4}}\right).
\]
The central order $m=\operatorname{ord}_{s=1/2}\zeta_N(s)$ is
positive and congruent to $2$ modulo $4$. For the shifted
Dirichlet-series deformation of $\zeta_N$, and every sufficiently
small positive $t$, its $m$ branches from the central zero consist
of exactly $(m+2)/2$ right branches and $(m-2)/2$ left branches.
The two initially vertical branches both bend right. This assertion
uses neither A nor M and does not assume $m=2$.
\end{corollary}

\begin{proof}
The field, its totally real tame $Q_8$ Galois group, ramification
only at $5,41$, and the nonfreeness of $\mathcal O_N$ over
$\mathbb Z[Q_8]$ are Martinet's Example~3, printed p.~407, with
Theorem~IV.1 on p.~406. They are classical inputs, already used for
central existence in prior MT-ART10; central existence is not the
new claim here. Fr\"ohlich's Theorem~1, printed p.~82 of his 1974
paper, identifies this nontrivial module class with root number
$-1$ for the unique irreducible character $\chi$ of degree two.

Choose a cyclic subgroup of order four. Its faithful linear
character induces $\chi$, so Artin induction identifies $L(s,\chi)$
with a nontrivial finite-order Hecke $L$-function over the quadratic
fixed field. It is entire. Its selfdual functional equation with
sign $-1$ makes its central order $\ell$ positive and odd. Its
Dirichlet coefficients are real, so its leading central coefficient
$b_\ell$ is a nonzero real number. The regular representation gives
the exact all-prime factorization
\[
 \zeta_N(s)=\zeta(s)L(s,\chi_5)L(s,\chi_{41})L(s,\chi_{205})
                     L(s,\chi)^2.
\]
The included 70-digit Arb certificate, using the finite Hurwitz
formula, proves
\[
 0.23<L(1/2,\chi_5)<0.24,\quad
 1.10<L(1/2,\chi_{41})<1.11,\quad
 1.17<L(1/2,\chi_{205})<1.18.
\]
Also $\zeta(1/2)<0$: the alternating eta series is positive by
pairing successive terms and $1-2^{1-s}<0$ for $0<s<1$.
It follows that $m=2\ell\equiv2\pmod4$ and the leading coefficient
of $\zeta_N$ is
$a_m=\zeta_F(1/2)b_\ell^2<0$, without a zero-location hypothesis.

The discriminant is $D_N=205^6$. For completeness, at either
$p=5$ or $41$ the inertia image in
$\operatorname{Gal}(F/\mathbb Q)=Q_8/\{\pm1\}$ is nontrivial.
Tame inertia in $Q_8$ is cyclic; any element with nontrivial image
in that quotient has order four. Thus the inertia order is four
and the tame discriminant exponent is $8(1-1/4)=6$. There are no
other ramified primes. This also agrees with Martinet's
Propositions~III.1--III.2, printed pp.~404--405.

The Euler product gives $\zeta_N'/\zeta_N(3/2)<0$ strictly.
Since $N$ is totally real of degree eight, the same certificate gives
\[
 C_N<2+3\log205+8a_{\mathbb R}
       <-3.5197<0.
\]
The theorem now gives the stated right/left branch counts.
\end{proof}

\paragraph{Source and verification locators.}
Martinet, \emph{Modules sur l'alg\`ebre du groupe quaternionien},
Ann. Sci. ENS~4 (1971), 399--408,
DOI~10.24033/asens.1216, Propositions~III.1--III.2,
Theorem~IV.1 and Example~3. The original p.~407 image was checked
again here: its second radical is $\sqrt{5\cdot41}$.
Fr\"ohlich, \emph{The Galois module structure of algebraic integer
rings in fields with generalised quaternion group},
M\'em. SMF~37 (1974), 81--86, DOI~10.24033/msmf.132,
Theorem~1 p.~82. This is an imported classical theorem; its
referenced 1972 proof is not claimed to have been re-proved here.
The short local deduction above is complete.
The script \path{check_central_jets.py} uses exact Kronecker
characters and
\[
 L(s,\chi_d)=d^{-s}\sum_{a=1}^d\chi_d(a)\zeta(s,a/d)
\]
with complex-ball Hurwitz evaluation, and writes every sign interval
to \path{central_jets_certificate.json}.

\paragraph{Geometric content and limitation.}
The central zero has local divisor length $m$; after the ramified
base change $t=u^m$, the normalization has $m$ sections. The
coefficient of $u^2$ depends on the completion unit and the shifted
Euler value, data not retained by the reduced support. On the chosen
positive real parameter ray the theorem assigns a canonical
multiset of sides. It does not assign a single unramified Boolean
label to the central support point, nor a colouring invariant under
complex monodromy. In the example, a nontrivial integral
$\mathbb Z[Q_8]$-module class supplies the odd Artin order, while
the analytic second jet supplies the excess of two right branches.
Those are precise distinct steps; no finite pure motive or
complexity lower bound is inferred.
